\section{A generalised constant Q transform}\label{sec_CSQT}

Before introducing frequency and time focus, we first introduce a slight generalisation of a standard constant Q transform (CQT for short). We consider the general framework of continuous time-frequency/time-scale transforms, the prototype of which is given by the Short Time Fourier Transform (STFT, see~\cite{grochenig2013foundations} for a detailed exposition) and the continuous wavelet transforms~\cite{Grossmann1984decomposition}.
%also belong to this general framework, which was given a unifying point of view in terms of group representations~\cite{Grossmann1985transforms,Grossmann1986transforms}.  
The constant $Q$ transform was introduced in a discrete context~\cite{Brown1991calculation} and revisited more recently~\cite{velasco2011constructing,holighaus2012framework}, we give below a slightly more general version adapted to the continuous setting.

The purpose of this section is to introduce the first tools used for our non linear transform but also some methods and computations that will be directly applied in section~\ref{sec_frequency_focus}.

\subsection{Definition of the time-frequency atoms}
The time-frequency atoms are built from a reference waveform $h$, which will be assumed continuous and compactly supported in the Fourier domain.  From this, time-frequency atoms $h_{t,\u}$ are generated as
\begin{equation}
	\forall x,t,\u\in\R\ ,\quad \ h_{t,u}(x) := \sqrt{\gamma(u)}e^{2i\pi\gamma(u)x}h(\gamma(u)x-t)\ ,
\end{equation}
where $\u$ is a frequency index, mapped to a frequency variable by the function $\gamma$.
%In this section and Section~\ref{sec:freqfocus} above,  the $u$ variable has to be seen as an index variable for the frequency which is at the end the values $\gamma(u)$.
\begin{remark}
    To make sense from  a physical point of view, the variable $t$ has to be seen as a dimensionless time variable, therefore $\gamma(\u)$ has to be understood as a dimensionless frequency.
\end{remark} 

The frequency function $\gamma$ is assumed to be a positive, increasing $C^1$ diffeomorphism, such that $\lim_{u\to -\infty}\gamma(u)=0$ and $\lim_{u\to +\infty}\gamma(u)=+\infty$. In~\cite{holighaus2012framework,velasco2011constructing}, $\gamma$ was given an exponential form, we consider here a slightly more general such scale function.

For technical reasons, the reference waveform $h$ will also be supposed to be symmetric\footnote{This assumption we be necessary for this section and the following one} with respect to $-1$ in the Fourier domain, which suggests to write $h$ as
\begin{equation}
\label{fo:hr}
    \forall x\in\R, \ h(x)= h_r(x)e^{-2i\pi x}\ ,
\end{equation}
where $h_r$ can be chosen real valued. 
%being a \textit{real-valued} window, assumed to be continuous, $L^2$ and compactly supported in the Fourier domain. 
An example of such a reference waveform could be a shifted version of frequency domain Hann, or raised cosine
\begin{equation}
	\hat{h}(\xi) = \mathds{1}_{[-3,1]}(\xi)\cos^2(\tfrac{\pi}{2} \xi)\ ,
\end{equation}
whose form in the time domain reads
\begin{equation}
	h(x) = \frac{\sin(4\pi x)}{4\pi x(1-4x^2)}e^{-2i\pi x}\ .
\end{equation}

We can now introduce the generalised Viennese constant Q transform~\cite{velasco2011constructing}. Let $f\in C_c(\R)$, then we define the transform through inner products with time-frequency atoms,
\begin{equation}
    \forall t,\u\in\R, \ Lf(t,\u) := \langle f , h_{t,\u} \rangle_{L^2_x}\ .
\end{equation}
%with the atoms defined by
%\begin{equation}
%	\forall x,t,\u\in\R, \ h_{t,\u}(x) := \sqrt{\gamma(\u)}e^{2i%\pi\gamma(\u)x}h(\gamma(\u)x-t)
%\end{equation}
The normalisation by $\sqrt{\gamma}$ ensures norm conservation, i.e. $\|h_{t,\u}\|_{L^2_x}=\|h\|_{L^2_x}$ for any $t,\u\in\R$. However such normalisation leads us to work in the modified $L^2(d\mu)$ space where the measure is defined by $d\mu(t,\u) := \tfrac{\gamma'(\u)}{\gamma(\u)}d\u dt$. The pointwise definition on $L^2$ is guaranteed by the fact that by assumption, $h$ is continuous and compactly supported in the frequency domain, therefore $h\in L^2$.


In order to determine the first norm control we need to compute the Fourier transform of the atoms and for the sake of simplicity let us introduce the auxiliary function $\beta_\xi$ defined by
\begin{equation}
\label{fo:def.beta}
\beta_\xi(\u):=\gamma(\u)^{-1}\xi-1\ .
\end{equation}
Now the Fourier transform of the atoms reads
\begin{equation}
	\F(h_{t,\u})(\xi) = \sqrt{\gamma(\u)}\int_\R e^{-2i\pi x (\xi-\gamma(\u))}h(\gamma(\u)x-t) dx\ .
\end{equation}
Using the change of variable $y=\gamma(\u)x-t$, we obtain
\begin{align}
	\F(h_{t,\u})(\xi) &= \frac{1}{\sqrt{\gamma(\u)}}\int_\R e^{-2i\pi \beta_\xi(\u)(y+t)}h(y)dy \\
	& = \frac{1}{\sqrt{\gamma(\u)}}e^{-2i\pi\beta_\xi(\u)t}\hat{h}(\beta_\xi(\u))\ .
\end{align}


\subsection{Admissibility condition and norm equality}

%Having the Fourier transforms of the atoms we can attack the first result of the journey which are the well-definedness of the operator from $L^2(\R,dx)$ into $L^2(\R^2,d\mu)$ and an explicit expression for the norm of the operator. 
Given these definitions and notations, we can now state the first results, namely the well-definedness of the operator from $L^2(\R,dx)$ into $L^2(\R^2,d\mu)$ and an explicit expression for the norm of the transform.
Such results are summed up in the following Proposition.
%
\begin{prop}\label{prop_norm_relation_operator_Q_transform}
Let $\gamma$ be a $C^1$ diffeomorphism such that $\lim_{\u\to -\infty}\gamma(\u)=0$ and $\lim_{\u\to +\infty}\gamma(\u)=+\infty$, and $h$ being a continuous and compactly supported in the Fourier domain. 
%and symmetric around $-1$ in Fourier\bruno{L'hypothèse de symétrie n'est pas nécessaire ici, ça peut faire une remarque plus loin}. 
Assume that $h$ satisfies the following admissibility conditions
\begin{equation}
    0<c_h := \int_{-1}^{+\infty}\frac{|\hat{h}(y)|^2}{y+1}dy = \int_{-\infty}^{-1}\frac{|\hat{h}(y)|^2}{-y-1}dy<\infty\ .
    \label{CQT.admissibility}
\end{equation}
Then for any $f\in L^2(\R,dx)$ we have
\begin{equation}
    \|Lf\|^2_{L^2(\R^2,d\mu)} = c_h \|f\|_{L^2(dx)}^2\ .
\end{equation}
\end{prop}
Notice that the right hand side equality in the above admissibility condition can be satisfied whenever $\hat h$ is symmetric with respect to -1, which is a simple and easy to fulfil condition.
\begin{proof}
	Let $f\in C_c^\infty$. Using Fubini's lemma we obtain 
	\begin{equation*}
		\|Lf\|_{L^2(d\mu)}^2 = \int\bar{\hat{f}}(\xi)\hat{f}(\xi')\hat{h}(\beta_\xi(\u))\bar{\hat{h}}(\beta_{\xi'}(\u))\frac{1}{\gamma(\u)}e^{-2i\pi(\gamma(\u)^{-1}(\xi-\xi'))t} d\xi d\xi' \frac{\gamma'(\u)}{\gamma(\u)}d\u dt\ .
	\end{equation*}
	By inverse Fourier transform this gives us that
	\begin{align*}
		\|L f\|_{L^2(d\mu)}^2 &= \int_\R|\hat{f}(\xi)|^2\int_\R \frac{\gamma'(\u)}{\gamma(\u)}|\hat{h}(\beta_\xi(\u))|^2d\u d\xi
	\end{align*}
     which concludes the first step of the proof.
	For the second part, let us set $y:=\beta_\xi(\u)$. Hence $\u=\gamma^{-1}(\frac{\xi}{y+1})$ and $d\u =-\frac{\gamma^2(\u)}{\xi\gamma'(\u)}dy$. Since the limits $\lim_{\u\to\pm\infty}\beta_\xi(\u)$ depend on the sign of $\xi$ we have to split the integral in two terms.
	\begin{align*}
		\|L f\|_{L^2}^2 &= \int_\R|\hat{f}(\xi)|^2\int_\R\frac{\gamma'(\u)}{\gamma(\u)}|\hat{h}(\beta_\xi(\u))|^2d\u d\xi \\
		& = \underset{=:(a)}{\underbrace{\int_{\R^+}|\hat{f}(\xi)|^2\int_\R\frac{\gamma'(\u)}{\gamma(\u)}|\hat{h}(\beta_\xi(\u))|^2d\u d\xi}} + \underset{=:(b)}{\underbrace{\int_{\R^-}|\hat{f}(\xi)|^2\int_\R\frac{\gamma'(\u)}{\gamma(\u)}|\hat{h}(\beta_\xi(\u))|^2d\u d\xi}}\ .
	\end{align*}
	For the first term $(a)$ we have $y(\u)=\frac{\xi}{\gamma(\u)}-1$ so $\lim_{u\to -\infty}y(u)=+\infty$ and $\lim_{u\to +\infty}y(u)=-1$ thus by the same change of variables as above we have
	\begin{align*}
		(a) &= \int_{\R^+}|\hat{f}(\xi)|^2\int_{-1}^{+\infty}\frac{\gamma'(\u)}{\gamma'(\u)}\frac{\gamma^2(\u)}{\xi\gamma(\u)}|\hat{h}(y)|^2dy \\
		&= \int_{\R^+} |\hat{f}(\xi)|^2 d\xi\int_{-1}^{+\infty}\frac{|\hat{h}(y)|^2}{y+1}dy = c_h \int_{\R^+} |\hat{f}(\xi)|^2 d\xi\ .
	\end{align*}
	The evaluation of the second term $(b)$ is similar but - due to the fact that $\lim_{u\to -\infty} y(u)=-\infty$ and $\lim_{u\to +\infty}=-1$ - we obtain at the end
	\begin{equation*}
		(b) = \int_{\R^-} |\hat{f}(\xi)|^2 d\xi\int_{-\infty}^{-1}\frac{-|\hat{h}(y)|^2}{y+1}dy = c_h \int_{\R^-} |\hat{f}(\xi)|^2 d\xi\ .
	\end{equation*}
	This gives us after the change of variables $y':=-2-y$ and by condition~\eqref{CQT.admissibility},
	\begin{equation*}
		(b) = \int_{\R^-} |\hat{f}(\xi)|^2 d\xi\int_{-1}^{+\infty}\frac{|\hat{h}(y)|^2}{y+1}dy\ ,
	\end{equation*}
	which allows us to conclude by summing up the two terms and by density of $C_c^\infty$ in $L^2$.
\end{proof}
    One purpose of such result is to guarantee the existence of a left inverse operator thanks to the classical Hilbertian analysis result, which we quote here for the sake of completeness.
\begin{prop}[Proposition 2.17 from \cite{conway2019course}]
\label{prop_inversion_isometry}
	Let $H,H'$ be $\K-$Hilbert spaces and $T:H\rightarrow H'$ and linear map. Then we have the existence of a left inverse if and only if $T$ is an isometry, furthermore such inverse is the conjugate map $T^*=\overline{T^t}$.
\end{prop}
Following this idea we can compute the proper inverse operator which is given as follows. Such result is similar to the one we can find in the classical time frequency transforms.
\begin{prop}[Inversion formula]\label{prop_inversion_formula_linear}
    The following inversion formula is true in $L^2$.
    \begin{equation}
f = \frac{1}{c_h}\int Lf(t,\u)h_{t,\u}\,d\mu(\u,t)\ .
\end{equation}

\end{prop}

\begin{remark}
\begin{enumerate}[(a)]
\item 
The admissibility condition~\eqref{CQT.admissibility} is reminiscent of the continuous wavelet transform (CWT). Indeed, expressed in terms of $h_r$ defined in~\eqref{fo:hr}, $Lf$ is actually a CWT of $f$ with wavelet $h_r$ as defined in~\eqref{fo:hr}, combined with a frequency scaling $\gamma$ that may differ from the usual $\gamma(\u)=2^\u$. The admissibility condition then becomes the usual wavelet admissibility condition on $h_r$.
\item 
Like the STFT and the continuous wavelet transform, alternative inversion formulas can be obtained by using different windows for the analysis and the synthesis, say $h$ as before and $w$. With the very same calculations as above, one obtains
\begin{equation*}
    f = \frac{1}{c_{hw}}\int Lf(t,\u)w_{t,\u}\,d\mu(t,\u)\ ,
    \end{equation*}with
\begin{equation*}
c_{hw} = \int_{-1}^\infty \frac{\overline{\hat h}(y)\hat w(y)}{y+1} dy= \int_{-\infty}^{-1} \frac{\overline{\hat h}(y)\hat w(y)}{-y-1} dy\ ,
\end{equation*}
assuming $c_{hw}$ is nonzero and the corresponding integrals converge.

\end{enumerate}
\end{remark}